\input{../tex-inputs/latex-korrekturansicht-vorspann}

\section[Arthur Schnitzler an Theodor Herzl, 19. 9. 1893]{L03905 Arthur Schnitzler an Theodor Herzl, 19. 9. 1893}
\nopagebreak\mylabel{L03905v}
\rehead{ }\normalsize\beginnumbering\briefempfaengerindex{Herzl, Theodor@\textsc{Herzl, Theodor}!zzzGoldmann, Paul@\emph{von Paul Goldmann}!1893-09-191@{19. 9. 1893}|(be}\briefempfaengerindex{Herzl, Theodor@\textsc{Herzl, Theodor}!zzzSchnitzler, Arthur@\emph{von Arthur Schnitzler}!1893-09-191@{19. 9. 1893}|(be}
\toendnotes[C]{\smallbreak\pagebreak[2]}
\correspDesc{Versand  durch Arthur Schnitzler, Paul Goldmann am 19. 9. 1893 in Salzburg
\newline{}Übermittlung  am 19. 9. 1893 in Salzburg
\newline{}Erhalt  durch Theodor Herzl im Zeitraum [20. 9. 1893 – 24. 9. 1893?] in Baden bei Wien}\toendnotes[C]{\smallbreak}
\Standort{Jerusalem, Central Zionist Archives, H1:1926-9.}
\physDesc{Brief, 3 Blätter, 4 Seiten, Kuvert, 1313 Zeichen
\newline{}Handschrift: schwarze Tinte, deutsche Kurrent
\newline{}Versand: Stempel: »\nobreak{}\oindex{Salzburg@\textbf{Salzburg}, \emph{Verwaltungsgebiet}|pwk}Salzburg-Stadt, 19 9 93, 12–\textcolor{gray}{1}M\nobreak{}«.  
\newline{}Ordnung: 1) mit Bleistift von unbekannter Hand 
                                 innerhalb das Konvoluts paginiert:
                                 »33«–»34«  2) mit Bleistift von unbekannter Hand bei der Datumsangabe ergänzt: »{[}1893{]}«}\toendnotes[C]{\smallbreak}\pstart{}{\pb}\textsc{Herrn Doctor Theodor Herzl}\pend{}\pstart{}\textcolor{pink}{\textsc{Baden bei} Wien}\oindex{Baden bei Wien@\textbf{Baden bei Wien}, \emph{Hauptstadt}|pw}{}\ledrightnote{\textcolor{pink}{Baden bei Wien}}\pend{}\pstart{}\textsc{\textcolor{pink}{Helenenstraße 71}\oindex{Helenenstraße 71@\textbf{Helenenstraße 71}, \emph{Wohngebäude}|pw}{}\ledrightnote{\textcolor{pink}{Helenenstraße 71}}.}\pend{}{\bigskip}\vspace{1em}
\pstart
           {\pb}\textcolor{pink}{Salzburg}\oindex{Salzburg@\textbf{Salzburg}, \emph{Verwaltungsgebiet}|pw}{}\ledrightnote{\textcolor{pink}{Salzburg}} den 19. September\pend
           
\pstart{}Mein verehrtes Freund,\pend\vspace{0.5em}
\pstart
           ich will Sie noch, bevor ich die Freude habe, Sie in \textcolor{pink}{Baden}\oindex{Baden bei Wien@\textbf{Baden bei Wien}, \emph{Hauptstadt}|pw}{}\ledrightnote{\textcolor{pink}{Baden bei Wien}} zu ſprechen, von hier aus aufs allerherzlichſte grüßen. Sie errathen
               bereits, daſs ich hier mit \textcolor{blue}{Paul Goldmann}\pwindex{Goldmann, Paul 31.\,1.\,1865 Breslau – 25.\,9.\,1935 Wien@\textsc{Goldmann, Paul} (31.\,1.\,1865 Breslau – 25.\,9.\,1935 Wien), \emph{Schriftsteller, Journalist}|pw}{}\ledrightnote{\textcolor{blue}{Paul Goldmann}} zusa{\geminationm}en
               getroffen bin, und zweifeln auch gewiſs nicht daran, daſs es vom Augen{\pb}blick meiner Ankunft an bis heute morgen, – Tag der
               Abreiſe – ununterbrochen geregnet hat. So konnten wir kaum aus der Stadt heraus, und
               haben in Kaffehäuſern und in den Straßen die letzten 2 oder 3 Jahre durchgeplaudert.
               Es war ſehr ſchön. Geſtern hab ich auch, die Intimität benützend, über Pauls Schulter
               weg, Ihren \label{K_L03905-1v}\edtext{Brief an ihn}{\lemma{\textnormal{\emph{Brief an ihn}}}\Cendnote{\textnormal{nicht überliefert}}}\label{K_L03905-1} gelesen – {\pb}und ka{\geminationn} nun nicht weiterſchreiben, weil man\footnote{\noindent{}Man = Paul Goldmann} neben mir ununterbrochen
               das \textsc{\textcolor{green}{Intermezzo}\pwindex{Verga, Giovanni 2.\,9.\,1840 Catania – 27.\,1.\,1922 ebd.@\textsc{Verga, Giovanni} (2.\,9.\,1840 Catania – 27.\,1.\,1922 ebd.), \emph{Schriftsteller}!Cavalleria rusticana. Oper in einem Aufzuge@\strich\emph{Cavalleria rusticana. Oper in einem Aufzuge}|pw}\pwindex{Targioni-Tozzetti, Giovanni 17.\,3.\,1863 Livorno – 30.\,5.\,1934 ebd.@\textsc{Targioni-Tozzetti, Giovanni} (17.\,3.\,1863 Livorno – 30.\,5.\,1934 ebd.), \emph{Librettist}!Cavalleria rusticana. Oper in einem Aufzuge@\strich\emph{Cavalleria rusticana. Oper in einem Aufzuge}|pw}\pwindex{Menasci, Guido 24.\,3.\,1867 Livorno – 27.\,12.\,1925 ebd.@\textsc{Menasci, Guido} (24.\,3.\,1867 Livorno – 27.\,12.\,1925 ebd.), \emph{Librettist}!Cavalleria rusticana. Oper in einem Aufzuge@\strich\emph{Cavalleria rusticana. Oper in einem Aufzuge}|pw}{}\ledrightnote{\textcolor{green}{Cavalleria rusticana. Oper in einem Aufzuge}}} pfeift und dreinredet.\pend
           
\pstart
           Auf baldiges Wiederſehen{\\[\baselineskip]}herzlich Ihr \spacefill\mbox{ArthSchnitzler}\pend
           \leftskip=0em{}\selectlanguage{ngerman}\vspace{1em}
\pstart{}{\pb}{[}hs. Goldmann:{]} Lieber
                  Freund!\pend\vspace{0.5em}
\pstart
           Ich danke Ihnen für Ihren ſo ſehr ſchönen Brief, auf den ich heute nicht antworte,
               weil ich heut zu dumm bin. Dies ſoll nur eine Empfangsbeſtätigung und ein Gruß ſein.
               Arthur Schnitzler hat das Verſprechen, Ihnen keinen geiſtreichen Brief zu ſchreiben,
               in niedrigſter Weiſe gebrochen. Ich kann alſo erſt recht nicht ſchreiben, weil ich
               nicht abſtechen mag. Glückliche Heimreiſe alſo! Geben Sie mir einen \textsc{\textcolor{gray}{conſt} de}{ }\textsc{téléphone} bei der Ankunft.\pend
           
\pstart
           Herzlichſt Ihr{\\[\baselineskip]}\spacefill\mbox{Paul Goldmann}\pend
           \leftskip=0em{}\selectlanguage{ngerman}\endnumbering\briefempfaengerindex{Herzl, Theodor@\textsc{Herzl, Theodor}!zzzGoldmann, Paul@\emph{von Paul Goldmann}!1893-09-191@{19. 9. 1893}|)be}\briefempfaengerindex{Herzl, Theodor@\textsc{Herzl, Theodor}!zzzSchnitzler, Arthur@\emph{von Arthur Schnitzler}!1893-09-191@{19. 9. 1893}|)be}\mylabel{L03905h}  \newcommand{\dateiname}{L03905}\newcommand{\titel}{Arthur Schnitzler an Theodor Herzl, 19. 9. 1893}\newcommand{\editorInnen}{Herausgegeben von Jahnke, SelmaMüller, Martin Anton}\input{../tex-inputs/latex-korrekturansicht-abspann}
